
%%%%%%%%%%%%%%%%%%%%%%%%%%%%%%%%%%%%
% Recap/Agenda
%%%%%%%%%%%%%%%%%%%%%%%%%%%%%%%%%%%%
% TODO better formatting
\begin{frame}
    \frametitle{Module 5: Beyond Recipes: Frequentist and Bayesian Approaches to Statistical Modeling}
    \begin{itemize}
        \item \hl{Previously: }Bayesian modeling with discrete distributions
        \item \hl{This time: }Bayesian modeling with the beta-binomial model
        \item \hl{Reading: }(none)
    \end{itemize}

\end{frame}
%%%%%%%%%%%%%%%%%%%%%%%%%%%%%%%%%%%%
% Learning objectives:
%%%%%%%%%%%%%%%%%%%%%%%%%%%%%%%%%%%%
\begin{frame}
    \frametitle{Learning Objectives}
    \begin{itemize}
        \item \textbf{M5, LO2: The Bayesian Perspective:} Explain how Bayesian modeling uses likelihoods, priors, and posteriors to perform inference. Explain how Bayesian analyses can be updated sequentially as new data arrives and how the results remain invariant to the order in which the data is observed.
        \item \textbf{M5, LO5: Bayesian Modeling 2:} Apply and interpret the beta-binomial model, including understanding the prior and posterior distributions, and evaluating the influence of the prior relative to the likelihood.
    \end{itemize}
\end{frame}
